\documentclass[a4paper,12pt]{article}
\usepackage[T1]{fontenc}
\usepackage[utf8]{inputenc}
\usepackage[polish]{babel}
\usepackage{lmodern}
\usepackage{geometry}
\usepackage{setspace}
\usepackage{enumitem}
\geometry{margin=2.5cm}
\setstretch{1.15}
\setlength{\parindent}{0pt}
\setlength{\parskip}{0.6em}

\title{Wykład: Pamięć w C++}
\author{Notatki do zajęć}
\date{\today}

\begin{document}

\maketitle

\section*{Wstęp}
Zarządzanie pamięcią to jeden z kluczowych tematów w języku C++.
W praktyce decyduje ono nie tylko o wydajności programu, ale również o jego
stabilności i bezpieczeństwie.

W C++ programista ma dużą kontrolę nad czasem życia obiektów oraz sposobem
alokacji pamięci. Ta elastyczność jest ogromną zaletą, ale wymaga dobrych
nawyków i znajomości mechanizmów języka.

\section*{Punkty wyróżnione}
\begin{itemize}[leftmargin=1.2cm]
    \item \textbf{AddressSanitizer (ASan)} -- narzędzie do wykrywania błędów
    związanych z pamięcią, m.in. wycieków, użycia pamięci po zwolnieniu
    (\textit{use-after-free}) oraz przepełnień bufora.
\end{itemize}

\end{document}
